\section{What a client can survive}
\label{sec:ch05-what-a-client-can-survive}
\index{schema!evolution|(}

\code{products} returns every row in the catalog. That was defensible at
twenty-five products and stops being defensible at some number nobody will
notice passing, and chapter~\ref{ch:data-without-the-n-plus-1} already built the
replacement. So the field has to go, and there is no version number in the URL
to hide the removal behind.

GraphQL schemas are versionless, which is true and is usually where the sentence
stops. The part that matters is what replaces the version: a GraphQL schema
gives its compatibility guarantees one field at a time, so somebody has to know,
for each change, which of three kinds it is.
Figure~\ref{fig:ch05-change-kinds} is that taxonomy, with this chapter's changes
placed on it.

\begin{figure}[htbp]
  \centering
  % Three kinds of schema change, with Mosaic's three placed on them. Referenced
% from chapter 05. Bare tikzpicture: the figure environment, caption and label
% live at the call site. Uses only arrows.meta and positioning.
\begin{tikzpicture}[
    font=\small,
    band/.style={draw, rounded corners=2pt, minimum width=34mm,
                 minimum height=11mm, align=center, font=\scriptsize},
    safe/.style={band},
    warn/.style={band, fill=black!8},
    bad/.style={band, fill=black!18, thick},
    ex/.style={font=\scriptsize\ttfamily, anchor=west, align=left},
    note/.style={font=\scriptsize\itshape, align=left, anchor=west},
    flow/.style={-{Stealth[length=2mm]}, thick}
  ]

  % ---- the three bands -------------------------------------------------
  \node[safe] (add)  at (0,0)    {additive\\nothing to do};
  \node[warn] (dep)  at (0,-1.9) {replaceable\\deprecate, then remove};
  \node[bad]  (brk)  at (0,-3.8) {breaking\\no way to soften it};

  \draw[flow] (add) -- (dep);
  \draw[flow] (dep) -- (brk);

  % ---- what each one means ---------------------------------------------
  \node[note, text width=52mm] at (2.2,0)
    {a new field, a new type,\\a new optional argument};

  \node[note, text width=52mm] at (2.2,-1.9)
    {the old thing still answers\\while callers move off it};

  \node[note, text width=52mm] at (2.2,-3.8)
    {a name, a type or a\\nullability that changed};

  % ---- Mosaic's three --------------------------------------------------
  \node[font=\scriptsize\bfseries, anchor=west] at (-2.6,-5.3) {in this chapter};

  \node[ex] at (-2.6,-6.0) {browseProducts};
  \node[note, text width=40mm] at (0.6,-6.0) {added last chapter, additive};

  \node[ex] at (-2.6,-6.7) {products};
  \node[note, text width=40mm] at (0.6,-6.7) {deprecated, still answers};

  \node[ex] at (-2.6,-7.4) {Product.id};
  \node[note, text width=40mm] at (0.6,-7.4) {same SDL, new format};

  \node[ex] at (-2.6,-8.1) {Product.reviews};
  \node[note, text width=40mm] at (0.6,-8.1) {list became a connection};

  % the two that bite
  \draw[flow, dashed] (-2.9,-7.4) -- (-3.6,-7.4) -- (-3.6,-3.8) -- (brk.west);
  \draw[flow, dashed] (-2.9,-8.1) -- (-4.0,-8.1) -- (-4.0,-4.2)
    -- (-1.7,-4.2) -- (brk.south west);

\end{tikzpicture}

  \caption{Three kinds of change, and where four of Mosaic's edits landed. The
  top band costs nothing and is where most work should live. The middle band is
  the only place a replacement can happen safely, and it is a process rather
  than an edit. The bottom band cannot be softened, only decided on, and two of
  this chapter's three changes are in it.}
  \label{fig:ch05-change-kinds}
\end{figure}

\index{schema!additive change}%
The top band is where a healthy schema spends most of its life. Adding a field,
a type, an enum value that nobody is switching on exhaustively, or an argument
with a default costs nothing, because a client asks for what it wants and is
handed exactly that. A field it has never heard of is a field it never
selects. \code{browseProducts} arrived in
chapter~\ref{ch:data-without-the-n-plus-1} this way and no client noticed.

The bottom band is anything that changes a name, a type, or a nullability in the
direction that gives the client less. There is no gentle version of it. A field
that was \code{[Review!]!} and is now \code{ReviewConnection!} does not degrade;
it fails.

\subsection*{Deprecating a field costs one attribute}

\index{deprecation!GraphQLDeprecated}%
The middle band is where \code{products} goes, and it is the reason the top and
bottom bands are not the whole story.

\begin{minted}[fontsize=\footnotesize]{csharp}
[GraphQLDeprecated(
    "Returns the whole catalog with no upper bound on its size. "
    + "Use `browseProducts`, which pages, filters and sorts.")]
public static Task<IReadOnlyList<Product>> GetProductsAsync(
    CatalogService catalog,
    CancellationToken cancellationToken)
    => catalog.GetProductsAsync(cancellationToken);
\end{minted}

That is the entire change. In the exported schema:

\begin{graphqlsdl}
products: [Product!]!
  @cost(weight: "10")
  @deprecated(
    reason: "Returns the whole catalog with no upper bound on its size. Use `browseProducts`, which pages, filters and sorts."
  )
\end{graphqlsdl}

\index{deprecation!behaviour}%
And the field goes on answering. That is worth stopping on, because the word
suggests otherwise. \code{@deprecated} turns nothing off. It is a message to the
people and the tools reading the schema, it shows up in introspection, and every
query that used the field yesterday works today. What has changed is that a
client generating code from this schema now gets a warning, and that the day the
field is removed becomes a day somebody can see coming.

Two smaller decisions are embedded in that attribute and both are worth copying.
The first is that the reason names the replacement. A deprecation that says
\enquote{do not use this} without saying what to use instead leaves every
consumer to work the answer out separately, and some of them will work it out
wrong. The second is that .NET's own \code{[Obsolete("reason")]} would have done
the same job: for fields, arguments, input fields and enum values the two
attributes behave identically~\autocite{chillicream2026versioning}. Mosaic uses
the GraphQL one because the deprecation is a fact about the schema rather than
about the C\#, and because the two stop being interchangeable the moment the
target is a type rather than a field.

\subsection*{What cannot be deprecated, and why the rule is right}

\index{deprecation!limits}%
There is one restriction, it is easy to trip over, and the reasoning behind it
is better than most: a non-null argument or input field with no default value
cannot be deprecated~\autocite{chillicream2026versioning}. The same restriction
applies to \code{@requiresOptIn}.

Think about what deprecating one would mean. The client is required to send it.
The server is telling the client to stop sending it. Both cannot be true, and
the day the argument is removed every query that supplies it stops parsing. So
the schema refuses to let you write down the contradiction. If you want to
retire a required argument, the move is to give it a default first, which is
additive, and deprecate it in the release after.

\index{deprecation!object types}%
Object types are the other gap, and this one is a live edge of the
specification rather than a settled rule. Deprecating a type itself needs an
option turned on in \code{ModifyOptions}, and it is off by
default because it is not in the released GraphQL specification at all: it
tracks an open proposal on the specification
repository~\autocite{graphqlspec2026objectdeprecation}. Mosaic leaves it off.
Two details are worth knowing before anybody turns it on. \code{[Obsolete]} on a
class deliberately does \emph{not} deprecate the type, only
\code{[GraphQLDeprecated]} does, because honouring \code{[Obsolete]} here too
would have silently deprecated types across every existing codebase the moment
the option was enabled. And a field that is not itself deprecated cannot return
a deprecated type; the schema fails to build. Reaching one through an interface
or a union is fine~\autocite{chillicream2026versioning}.

\subsection*{The other direction}

\index{schema!requiresOptIn}%
\code{@deprecated} marks what is leaving. There is a matching directive for what
has not arrived yet, and it gets far less attention than it deserves.
\code{@requiresOptIn} marks a field as unstable, hides it from introspection,
and requires a consumer to name the feature before they can see it. Like object
deprecation it is off until you ask for it, behind its own schema
option~\autocite{chillicream2026versioning}. A schema can then carry an
experiment
without every client discovering it, generating a type for it, and treating it
as a promise.

Mosaic does not use it, which is a scoping decision rather than a judgement:
nothing here is experimental enough to need it. I am pointing at it because the
usual complaint about versionless schemas is that there is nowhere to put a
half-finished idea, and there is.

\subsection*{Nothing here is enforced}

\index{deprecation!enforcement}%
One admission before moving on, because the rest of this section reads more
confidently than the situation deserves. A deprecation in Mosaic is a string in
the schema. Nothing rejects a query that uses \code{products}, nothing counts
how many clients still do, and nothing stops the field being deleted tomorrow
by somebody who did not read the reason. Deprecation is only half a mechanism
until a registry is watching the schema and a check is failing a pull request
that removes a field somebody is still selecting.
Chapter~\ref{ch:ci-cd-for-a-federated-graph} is where that half arrives.

Which brings us to the first change in this chapter that a client cannot
survive, and the reason it is worth taking anyway.
\index{schema!evolution|)}
